\chapter*{Abstract}

Losses after harvest have always been a significant issue in developing states that lowers the income
of farmers and a restriction of food availability. This is more so in smallholder agribusinesses as a
result of poor storage conditions, low practices in handling, and the lack of access to affordable
technologies. The goal of this paper is to come up with a machine learning-based model of predicting
the risk of post-harvest loss based on low-cost and the availability of data.

The research design was the quantitative and experimental study based on the structured information
used in the Nigeria General Household Survey ~\cite{ghspanel2024nigeria}. The use of important
variables in it were the type of crop, manner of storage, length of storage, geographical area and
the alleged causes of loss. Three machine learning models which include a Logistic Regression model,
a Decision Tree, and a Random Forest were implemented and tested on the basis of accuracy, precision,
recall, F1-score, and AUC-ROC values.

The results show that the most important predictors of post-harvest losses are storage method and
storage duration, with the type of crop and location coming next. Random Forest was the best
predictive model and showed good capability to fit intricate data associations. Decision tree models
were also feasible as well as interpretable, while Logistic Regression served as a strong baseline.
As was shown in this paper, even the simplest models of machine learning can be applied to make good
predictions on whether there would be losses after harvest and issue correct tips to smallholder
farmers. The projections will be meant to provide an excellent guide toward improving decision making
in the storage, handling and marketing practices. The results of the study demonstrate the feasibility
of low-cost decision support systems and the possibility of creating data-driven solutions to improve
productivity, minimise losses, and promote food security in developing countries.
